% !PRINTED 439
\weberpage{439}
%<<CONT
hence
\[ 4 p^\varepsilon = y^2 + m x^2. \]

From this $x$ and $y$ are determined up to sign, and $x$ may be assumed positive, and the sign of $y$ follows from the congruence
\weq{7}{y \equiv b x \pmod{p^\varepsilon}.}

If, on the other hand, one considers the composition $l^{-1} k$, then one has to compose the form $(a, b, c p^\varepsilon)$ with $(p, -b, a c p^{\varepsilon - 1})$; one determines $b'$ from the two congruences:
\[ b' \equiv b \pmod{2a}, \qquad b' \equiv -b \pmod{2 p^\varepsilon}, \]
and everything is to be carried out in just the same way, except that in place of (7) the congruence % [sic] German print has the grammatically anomalous dative "der Kongruenz"
\weq{8}{y \equiv -b x \pmod{p^\varepsilon}}
enters, that is, $y$ receives the opposite sign.

But according to (5) and (6)
\weq{9}{\alpha + \beta\omega = \frac{y + (2 a \omega + b) x}{2} = \frac{y + \sqrt{-m}\, x}{2},}
in which, if $a$ and $x$ and consequently $\beta$ are assumed positive, $\sqrt{-m}$ is to be taken positive imaginary.

But by \S\,38, (4)
\[ \left( \frac{\eta(\omega_\varepsilon)}{\eta(\omega)} \right)^{24} = \left( \frac{y + \sqrt{-m}\, x}{2} \right)^{12}, \]
and consequently one obtains by (3)
\weq{10}{\Phi(k, lk) = \left( \frac{y + \sqrt{-m}\, x}{2} \right)^{12}.}

In the same way there follows, on account of (16): % [sic] printed (16); no antecedent equation (16) occurs in this section here
\weq{11}{\Phi(k, l^{-1}k) = \left( \frac{y - \sqrt{-m}\, x}{2} \right)^{12},}
and equation (10) is therefore \gesperrt{not satisfied} when $(lk)$ is replaced by $(l^{-1}k)$, except when
\weq{12}{\left( \frac{y + \sqrt{-m}\, x}{2} \right)^{12} = \left( \frac{y - \sqrt{-m}\, x}{2} \right)^{12}}
holds.

But this is not possible if, as we have assumed, $p$ does not divide $m$. For if one decomposes $p$ in the field $\Omega = \Re(\sqrt{D})$ into the two conjugate prime ideals $\mathfrak{p}, \mathfrak{p}'$, distinct from one another, then the one divides $\frac{1}{2}(y + \sqrt{-m}\, x)$, the other
%CONT>>
